\section*{Capítulo \ref{ch:b3:planetary-interiors-surfaces} --- Interiores e superfícies planetárias}

\begin{solution}{exo:b3:planetary-interiors-surfaces:1}
$\tfrac52C\bar\rho = 2.5\times0.35\times5429 = 4750$; $x^2 = (4750 -
3300)/(5429 - 3300) = 0.681$, $x = 0.825$, um núcleo de $0.825\times2440 =
\qty{2010}{km}$; $\rho_c = 3300 + 2129/0.825^3 = \qty{7090}{kg.m^{-3}}$.
\end{solution}

\begin{solution}{exo:b3:planetary-interiors-surfaces:2}
$m = 3500\times\tfrac43\pi(25)^3 = \qty{2.29e8}{kg}$; $E = \tfrac12 m v^2 =
\tfrac12\times\num{2.29e8}\times(\num{2.0e4})^2 = \qty{4.6e16}{J}$, ou
$\num{4.6e16}/\num{4.2e15} = \qty{11}{Mt}$. A regra dá uma cratera de
$20\times\qty{50}{m} = \qty{1}{km}$.
\end{solution}

\begin{solution}{exo:b3:planetary-interiors-surfaces:3}
$N_1 = 120/\num{100000} = \qty{1.2e-3}{km^{-2}}$, erro relativo $1/\sqrt{120}
= \qty{9}{\percent}$. Na parte linear, $T = N_1/c = \num{1.2e-3}/\num{4.875e-4}
= \qty{2.46}{Ga}$ (o termo exponencial, $\num{8e-7}$, é desprezível), e
$\sigma_T = T/\sqrt{120} = \qty{0.22}{Ga}$.
\end{solution}

\begin{solution}{exo:b3:planetary-interiors-surfaces:4}
Área $4\pi(\num{1.8215e6})^2 = \qty{4.17e13}{m^2}$; potência $2.5\times
\num{4.17e13} = \qty{1.04e14}{W} = \qty{104}{TW}$, 2.2 vezes o total da Terra
a partir de uma superfície doze vezes menor, e $2.5/0.092 = 27$ vezes o seu
fluxo.
\end{solution}

\begin{solution}{exo:b3:planetary-interiors-surfaces:5}
(a)~$\bar\rho = \num{6.417e23}/(\tfrac43\pi(\num{3.3895e6})^3) =
\qty{3934}{kg.m^{-3}}$; $g = GM/R^2 = \qty{3.73}{m.s^{-2}}$; $\rho gR =
\qty{4.97e10}{Pa}$. (b)~$19\mu/(2\rho gR) = 19\times\num{6.7e10}/\num{9.94e10}
= 12.8$, $k_2 = 1.5/13.8 = 0.109$. (c)~O valor medido é maior em
2.6 desvios-padrão: Marte responde mais do que um planeta sólido com a
rigidez do manto, como ocorreria se houvesse um núcleo líquido em seu interior.
\end{solution}

\begin{solution}{exo:b3:planetary-interiors-surfaces:6}
(a)~Lua: $4\times9.8/1.62 = \qty{24}{km}$ em rocha cristalina,
\qty{12}{km} em sedimentos; Marte e Mercúrio: $4\times9.8/3.7 =
\qty{10.6}{km}$ e \qty{5.3}{km}. Os \qty{20}{km} da Lua ficam entre
suas duas previsões: a regra vale até algumas dezenas por cento. (b)~$Y =
\rho gD/3 = 2600\times9.8\times4000/3 = \qty{34}{MPa}$, a resistência de
rocha fraturada, e não de um cristal intacto.
\end{solution}

\begin{solution}{exo:b3:planetary-interiors-surfaces:7}
Marte: $\num{3.22e-14}(e^{6.93T} - 1) + \num{4.875e-4}\,T = \num{2.0e-3}$;
a estimativa linear $4.1$~Ga é muito alta, sendo o termo exponencial grande
ali; a bissecção dá $T = \qty{3.34}{Ga}$ (termo exponencial
$\num{3.6e-4}$, termo linear $\num{1.63e-3}$). A função lunar dá
\qty{2.39}{Ga}. Cada função é calibrada para um corpo: a taxa de
impactos e o tamanho das crateras que um dado projétil forma diferem
entre a Lua e Marte.
\end{solution}

\begin{solution}{exo:b3:planetary-interiors-surfaces:8}
(a)~$3.551181/1.769138 = 2.0073$ e $7.154553/3.551181 = 2.0147$, próximos de
$2$. (b)~$0.565247 - 0.844789 + 0.279542 = 0.000000$ por dia, com
$\num{1e-7}$: os três movimentos médios estão travados, uma ressonância que mantém
as posições relativas das luas em cada conjunção. (c)~As forças de atração se repetem nos
mesmos pontos da órbita, de modo que se somam em vez de se anularem;
elas alimentam a excentricidade tão rapidamente quanto as marés a amortecem.
\end{solution}

\begin{solution}{exo:b3:planetary-interiors-surfaces:9}
$n = 2\pi/(1.769138\times86400) = \qty{4.111e-5}{s^{-1}}$, $nR =
\qty{74.9}{m.s^{-1}}$, $(nR)^5/G = \qty{3.53e19}{W}$;
$\tfrac{21}{2}e^2(nR)^5/G = 10.5\times\num{1.6e-5}\times\num{3.53e19} =
\qty{5.9e15}{W}$ por unidade de $k_2/Q$. Então $k_2/Q =
\num{1.04e14}/\num{5.9e15} = 0.018$.
\end{solution}

\begin{solution}{exo:b3:planetary-interiors-surfaces:10}
(a)~$x^2 = (\tfrac52C\bar\rho - \rho_m)/(\bar\rho - \rho_m)$ é linear em
$C$, com coeficiente angular $\tfrac52\bar\rho/(\bar\rho - \rho_m)$. (b)~O coeficiente angular é
$2.5\times3340/52 = 161$: $x^2$ varia $0.161$ para cada $0.001$ de $C$.
Para $C = 0.395$, $x^2 = 0.036 + 0.161 = 0.197$, $x = 0.44$, $\rho_c =
\qty{3880}{kg.m^{-3}}$; para $C = 0.393$, $x^2 < 0$: não há núcleo algum.
(c)~Perto de $C = \tfrac25$ o momento de inércia praticamente não sente um núcleo pequeno:
seu terceiro algarismo significativo decide entre um núcleo grande e leve e nenhum.
São necessárias ondas sísmicas, marés e o campo magnético.
\end{solution}

\begin{solution}{exo:b3:planetary-interiors-surfaces:11}
(a)~$\sigma_T = N_1/(10\,N_1')$: em \qty{1.0}{Ga}, $N_1 \approx N_1' \approx
\num{4.875e-4}$, $\sigma_T = \qty{0.10}{Ga}$; em \qty{3.8}{Ga}, $N_1 =
\num{1.07e-2}$ e $N_1' = \num{6.15e-2}$, $\sigma_T = \qty{0.017}{Ga}$,
seis vezes menor. (b)~Com $N_1 = cT$, $N_1' = c$ e $\sigma_T =
cT/(\sqrt n\,c) = T/\sqrt n$. (c)~A própria função de cronologia é
incerta: ela se apoia em poucos sítios datados da Lua e na extrapolação
da Lua para Marte, erros que nenhuma contagem reduz.
\end{solution}

\begin{solution}{exo:b3:planetary-interiors-surfaces:12}
(a)~A razão é $(n_ER_E/n_IR_I)^5(e_E/e_I)^2 = (0.4982\times0.8569)^5
\times(0.009/0.004)^2 = 0.0142\times5.06 = 0.072$: cerca de
\qty{7.5e12}{W}, ou $\num{7.5e12}/(4\pi(\num{1.5608e6})^2) =
\qty{0.24}{W.m^{-2}}$. (b)~Um oceano desacopla a casca de gelo do
interior: a casca se deforma muito mais do que uma lua sólida, e seu gelo
dissipa energia de forma diferente da rocha quente de Io, e tanto $k_2$ quanto $Q$ mudam.
\end{solution}

\begin{solution}{pb:b3:planetary-interiors-surfaces:1}
\textbf{1.} $\bar\rho = \num{6.4169e23}/(\tfrac43\pi(\num{3.3895e6})^3) =
\qty{3934}{kg.m^{-3}}$.

\textbf{2.} $\tfrac52C\bar\rho = 3600$; $x^2 = (3600 - 3500)/(3934 - 3500) =
0.230$, $x = 0.479$, um núcleo de \qty{1620}{km}; $\rho_c = 3500 +
434/0.479^3 = \qty{7450}{kg.m^{-3}}$.

\textbf{3.} $x = 1830/3389.5 = 0.540$; $\rho_m = (3600 - 0.2915\times
3934)/(1 - 0.2915) = \qty{3460}{kg.m^{-3}}$, $\rho_c = 3462 +
472/0.1574 = \qty{6460}{kg.m^{-3}}$.

\textbf{4.} É um terço menor. A pressão no núcleo de Marte é muito
menor do que na Terra, então seu metal é menos comprimido; e um núcleo tão
leve sugere elementos muito mais leves misturados ao ferro.

\textbf{5.} O \omterm{def:b3:planetary-interiors-surfaces:love-number}{número de Love} (2003) surgiu muito antes do InSight (2021): ele
já mostrava que pelo menos o núcleo externo é líquido e situava seu
raio entre 1520 e \qty{1840}{km}, um intervalo que contém o
valor sismológico.

\textbf{6.} $1 - (1 - 20/3389.5)^3 = \qty{1.8}{\percent}$ a $1 - (1 -
37/3389.5)^3 = \qty{3.2}{\percent}$ do volume, e menos da massa:
mais leve que o manto, reduz um pouco a densidade efetiva do manto,
um efeito menor que a incerteza do modelo.

\textbf{7.} A partir de \qty{1}{km}: 54; a partir de 1.41: 16; a partir de 2: 10; a partir de 2.83: 5;
a partir de 4, 5.66, 8 e 11.3: 2; a partir de 16: 1.

\textbf{8.} $N = 54/80000 = \qty{6.75e-4}{km^{-2}}$, $\sigma = \sqrt{54}/80000
= \qty{0.92e-4}{km^{-2}}$ (\qty{14}{\percent}).

\textbf{9.} Esperam-se $54\times0.355 = 19$ crateras de \qty{1.41}{km} e
$54\times0.147 = 7.9$ de \qty{2}{km}; contadas 16 e 10, dentro de um
desvio-padrão ($\sqrt{19} = 4.4$, $\sqrt{7.9} = 2.8$): consistente.

\textbf{10.} Não: $0.072/\num{6.75e-4} = 107$ vezes abaixo dela. A contagem
registra o tempo.

\textbf{11.} Intervalos de largura igual em $\log D$ cobrem a faixa de forma uniforme e
mantêm números comparáveis de crateras nos intervalos de uma distribuição íngreme.

\textbf{12.} São crateras secundárias, formadas por blocos ejetados de uma
grande cratera primária: não são impactos independentes vindos do espaço; eles acrescentariam
crateras que a função de cronologia não conta.

\textbf{13.} $T = \num{6.75e-4}/\num{4.875e-4} = \qty{1.385}{Ga}$; o
termo exponencial, $\num{3.22e-14}(e^{9.6} - 1) = \num{4.7e-10}$, é uma
milionésima parte de $N_1$.

\textbf{14.} $\sigma_T = T/\sqrt{54} = \qty{0.19}{Ga}$: $1.38 \pm 0.19$~Ga.

\textbf{15.} $\num{6.75e-4}/\num{8.38e-4} = \qty{0.81}{Ga}$, mais jovem por
quase a metade.

\textbf{16.} $31 + 12 + 7 + 2 + 0 + 2 + 2 = 56$ crateras; $N(\qty{2}{km}) =
56/120000 = \qty{4.67e-4}{km^{-2}}$; $N_1 = \num{4.67e-4}/0.1473 =
\qty{3.17e-3}{km^{-2}}$. Bissecção em $\num{3.22e-14}(e^{6.93T} - 1) +
\num{4.875e-4}\,T = \num{3.17e-3}$ dá $T = \qty{3.54}{Ga}$ (termo exponencial
\num{1.45e-3}, termo linear \num{1.73e-3}).

\textbf{17.} $N_1' = \num{3.22e-14}\times6.93\,e^{6.93\times3.54} +
\num{4.875e-4} = \num{1.05e-2}$; $\sigma_T = \num{3.17e-3}/(\sqrt{56}\times
\num{1.05e-2}) = \qty{0.040}{Ga}$.

\textbf{18.} Planície A: \qty{14}{\percent}; planície B: \qty{1.1}{\percent}. Na
parte íngreme e inicial da cronologia, uma grande variação da densidade de
crateras significa uma pequena variação da idade; na parte linear a idade é tão
incerta quanto a contagem.

\textbf{19.} $120000\times\num{3.17e-3}\times(1 - 0.1473) = 324$; ela
carrega $48 + 8 = 56$.

\textbf{20.} $N_1 = 56/(120000\times0.8527) = \qty{5.47e-4}{km^{-2}}$,
$T_r = \num{5.47e-4}/\num{4.875e-4} = \qty{1.12}{Ga}$.

\textbf{21.} $\sigma = 1.12/\sqrt{56} = \qty{0.15}{Ga}$.

\textbf{22.} Uma cratera de \qty{2}{km} tem \qty{400}{m} de profundidade: a lava tinha
no mínimo essa espessura. Crateras maiores são mais profundas e suas bordas
ficam mais altas; elas foram parcialmente preenchidas e ainda aparecem.

\textbf{23.} $1.38 \pm 0.19$ e $1.12 \pm 0.15$~Ga diferem por \qty{0.26}{Ga},
cerca de um desvio-padrão combinado ($\sqrt{0.19^2 + 0.15^2} = 0.24$):
as contagens permitem um único episódio, embora não o comprovem.

\textbf{24.} A planície de lava jovem A tem uma idade-modelo de
\textbf{1.4~bilhão~de~anos} ($1.38 \pm 0.19$~Ga).
\end{solution}
